\begin{thebibliography}{9}

\bibitem{blackard1993} K. L. Blackard, T. S. Rappaport y C. W. Bostian, ``Measurements and models of radio frequency impulsive noise for indoor wireless communications,'' \emph{IEEE J. Sel. Areas Commun.}, vol. 11, no. 7, pp. 991--1001, 1993.

\bibitem{middleton1977} D. Middleton, ``Statistical-physical models of electromagnetic interference,'' \emph{IEEE Trans. Electromagn. Compat.}, vol. EMC-19, no. 3, pp. 106--127, 1977.

\bibitem{zhan2022} M. Zhan, Z. Pang, D. Dzung, K. Yu y M. Xiao, ``Short-packet interleaver against impulse interference in practical industrial environments,'' arXiv:2203.00770, 2022.

\bibitem{proakis2008} J. G. Proakis y M. Salehi, \emph{Digital Communications}, 5.ª ed. Nueva York, NY, EE.\,UU.: McGraw-Hill, 2008.

\bibitem{ortega2019} H. Ortega Boada y Ó. M. Reyes Torres, \emph{Comunicaciones digitales basadas en radio definida por software}, Editorial UIS, 2019.

\bibitem{hackrf} Great Scott Gadgets, ``HackRF One documentation.'' [En línea]. Disponible: \url{https://hackrf.readthedocs.io/en/stable/hackrf_one.html}

\bibitem{uhd} Ettus Research, ``UHD manual.'' [En línea]. Disponible: \url{https://files.ettus.com/manual/}

\bibitem{gnuradio} GNU Radio Project, ``GNU Radio.'' [En línea]. Disponible: \url{https://www.gnuradio.org}

\end{thebibliography}
